

\chapter*{ABSTRACT}
\phantomsection
\addcontentsline{toc}{chapter}{ABSTRACT}


\bigskip

The proliferation of high-resolution image sensors in surveillance and monitoring applications has exposed a fundamental limitation in conventional edge inference pipelines: the spatial downsampling required to fit high-resolution frames into fixed-dimension neural network inputs systematically destroys the fine-grained detail necessary for reliable detection of small or distant subjects. This work presents a distributed tile-based inference architecture for real-time \gls{hpd} on commodity microcontroller hardware, designed to preserve spatial resolution while remaining within the memory and compute constraints of resource-limited edge nodes.
End-to-end benchmarking reveals that frame capture and tiling completes in 4\gls{ms}, aggregator classification in 0.342\gls{ms}, and feature transmission per tile in 472ms. The dominant bottleneck is tile-level inference on the ESP32-S3, which requires 2,657ms per tile due to the mandatory use of OCTAL PSRAM for the 451KB TFLite arena, yielding a total critical-path latency of approximately 18.9 seconds against a 1,000 \gls{ms} target frame budget. All three deployed models were verified to be fully INT8 with zero float32 tensors remaining. The results confirm architectural correctness across all pipeline stages while clearly identifying PSRAM-bound inference latency and sequential tile fragmentation as the primary targets for optimisation in future work.
\vspace{1em}

\noindent\textbf{Keywords:}\textit{Human presence detection, distributed tile-based inference, local binary convolutional neural network, attention mechanism, TinyML, post-training quantization, ESP32, edge computing.}
